\documentclass[12pt]{article}

\usepackage[utf8]{inputenc}
\usepackage[margin=1in]{geometry}
\usepackage{amsmath,amssymb,amsthm,mathtools}
\usepackage{booktabs}
\usepackage{threeparttable}
\usepackage{graphicx}
\usepackage{xcolor}
\usepackage[authoryear,round]{natbib}
\usepackage{setspace}
\usepackage[nopatch=footnote]{microtype}
\usepackage{caption}
\usepackage{subcaption}
\usepackage{float}

\definecolor{darkblue}{rgb}{0,0,0.55}

\usepackage[hyphens]{url}
\usepackage{hyperref}
\hypersetup{
    colorlinks=true,
    linkcolor=darkblue,
    citecolor=darkblue,
    urlcolor=darkblue
}

\onehalfspacing
\captionsetup{font={small}, labelfont={bf}, skip=8pt}

\newtheorem{theorem}{Theorem}
\newtheorem{lemma}{Lemma}
\newtheorem{proposition}{Proposition}
\newtheorem{corollary}{Corollary}

\title{\textbf{The Cognitive Rent Disruption: Generative AI Exposure, Task Displacements, and the Reversal of Skill-Biased Technical Change Across Regional Labor Markets}\thanks{I am grateful to colleagues and seminar participants for helpful discussions and comments. All empirical calculations utilize publicly available administrative records from the U.S. Bureau of Labor Statistics, the Current Population Survey, and exposure metrics from \citet{felten2021, felten2023}. All remaining errors are my own.}}

\author{
    \textbf{Danil Popov}\thanks{Independent Researcher. Correspondence: \href{mailto:perque.balom@gmail.com}{\texttt{perque.balom@gmail.com}}.} \\
    \textit{Independent Researcher}
}

\date{October 27, 2026}

\begin{document}

\maketitle

\begin{abstract}
\noindent For over four decades, the economic literature on technical change and wage inequality has been anchored by the hypothesis of Skill-Biased Technical Change (SBTC), which posits that new technologies complement highly educated cognitive workers while substituting for routine manual and clerical labor. This paper examines how the commercial diffusion of frontier Generative Artificial Intelligence (GenAI) and Large Language Models (LLMs) departs from this historical regularity. Unlike earlier waves of computerization and industrial robotics, foundation models automate non-routine symbolic, linguistic, and analytical synthesis that previously generated scarcity rents for college-educated professionals. I develop a general equilibrium task continuum model showing that when algorithmic capital encroaches on intermediate cognitive task bands, the college wage premium contracts and upper-tail within-occupation wage dispersion narrows faster than lower-tail dispersion. To test these predictions, I construct a balanced panel tracking 753 detailed 6-digit SOC occupations and 3,218 U.S. counties from 2018 to 2025 ($N = 6,024$ occupation-years and $N = 25,744$ county-years), linking cohort-deconvolved BLS OEWS and QCEW earnings records to Language Modeling Exposure (LME) and AI Geographic Exposure (AIGE) benchmarks. Using two-way fixed effects, triple-difference, and event-study designs around the late-2022 release of commercial LLMs, I find that post-2023, a one-standard-deviation increase in language modeling exposure is associated with a $1.32\%$ annual deceleration in occupational mean wages ($t = -5.28, p < 0.001$). Triple-difference estimates show that this compression is driven by college-credentialed occupations ($\hat{\beta}_{\text{DDD}} = -0.0148, p < 0.001$), whereas manual trades and in-person services exhibit wage resilience. Within college-educated occupations, upper-tail wage dispersion ($\ln(P90/P50)$) contracts by $2.45$ log points per standard deviation of exposure ($p < 0.001$), compared to a statistically insignificant $-0.42$ log points in the lower tail ($\ln(P50/P10)$). Across regional labor markets, employment-weighted wage growth in high-exposure metropolitan counties slowed by $0.86\%$ per standard deviation of exposure ($t = -5.73, p < 0.001$), initiating a measurable phase of spatial wage convergence.

\vspace{1em}
\noindent \textbf{Keywords:} Generative AI, Skill-Biased Technical Change, Cognitive Scarcity Rents, Wage Inequality, Task-Based Models, Spatial Convergence.

\vspace{0.5em}
\noindent \textbf{JEL Classification Codes:} J24, J31, O33, R11, R12.
\end{abstract}

\newpage

\section{Introduction}
\label{sec:introduction}

For the past half-century, empirical and theoretical research on technological progress and the wage structure has been organized around the canonical framework of Skill-Biased Technical Change (SBTC) \citep{tinbergen1974, katzMurphy1992, acemoglu1998, autor2003, goldinKatz2008, acemogluAutor2011}. In its standard formulation, technological innovation complements highly educated workers performing abstract, non-routine analytical tasks while substituting for middle- and lower-skill workers engaged in routine codifiable tasks. This mechanism has been widely shown to account for the secular rise in the college wage premium, the widening of upper-tail wage inequality, and the growing earnings divergence between knowledge-intensive metropolitan areas and the rest of the United States \citep{autor2008, moretti2012, acemogluRestrepo2022}.

In this paper, I present theoretical and empirical evidence that the diffusion of frontier Generative Artificial Intelligence (GenAI) and Large Language Models (LLMs) marks a structural departure from this pattern. Unlike industrial robotics, enterprise software, or traditional factory automation, foundation models operate directly on unstructured text, code, and symbolic logic. They automate non-routine cognitive tasks---including software engineering, contract drafting, financial modeling, paralegal research, and technical synthesis---that have historically constituted the comparative advantage of college-educated professionals. By contrast, tasks requiring physical dexterity in unstructured environments, on-site mechanical troubleshooting, or high-touch interpersonal care (such as electrical work, plumbing, construction trades, and clinical nursing) remain largely insulated from language-based automation due to Moravec's paradox. Consequently, generative models exert direct downward pressure on the scarcity rents traditionally captured by cognitive labor.

I make three contributions to the analysis of this technological transition. First, I formalize cognitive task encroachment within a general equilibrium task continuum framework building on \citet{acemogluAutor2011} and \citet{acemogluRestrepo2018, acemogluRestrepo2022}. Whereas canonical automation models assume that capital displaces labor monotonically from lower-skill routine intervals upward, Generative AI introduces displacement within an interior cognitive band $[x_H^0, x_H^1]$ located inside the high-skill task domain. I derive exact conditions under which an expansion of generative AI capabilities reduces the equilibrium college-to-non-college wage ratio $w_H / w_L$ (Proposition \ref{prop:inversion}). I then extend the model to incorporate within-occupation ability heterogeneity, showing that when foundation models establish a high algorithmic floor for intermediate cognitive synthesis while substituting for scarce human drafting speed, within-occupation wage dispersion compresses substantially faster in the upper tail ($\ln(P90/P50)$) than in the lower tail ($\ln(P50/P10)$) (Proposition \ref{prop:dispersion}).

Second, I test these predictions using a balanced longitudinal panel of 753 detailed 6-digit Standard Occupational Classification (SOC) occupations across the United States spanning May 2018 through May 2025 ($N = 6,024$ occupation-year observations). Combining cohort-deconvolved administrative records from the U.S. Bureau of Labor Statistics (BLS) Occupational Employment and Wage Statistics (OEWS) program and point-in-time validation from Current Population Survey Outgoing Rotation Groups (CPS-ORG) with the AI exposure benchmarks of \citet{felten2021, felten2023}, I separate general machine-learning exposure from specific Language Modeling Exposure (LME). Because commercial foundation models (beginning with ChatGPT in November 2022 and GPT-4 in March 2023) diffused after the May 2022 OEWS survey reference window, the 2018--2025 sample provides five pre-treatment waves (2018--2022) and three full post-treatment waves (2023--2025).

Estimating two-way fixed effects regressions, I find that a one-standard-deviation increase in Language Modeling Exposure is associated with a $1.32\%$ annual deceleration in occupational mean wages during the post-2023 period ($t = -5.28, p < 0.001$). Dynamic event-study specifications benchmarked to May 2022 demonstrate that this represents a sharp reversal of prior trends: between 2018 and 2021, occupations with higher language exposure enjoyed significant positive wage growth premiums ($+1.68\%$ per standard deviation in 2018 relative to 2022, $t = 7.64$), consistent with ongoing SBTC. Following the commercial deployment of LLMs in late 2022, this gradient turned negative, reaching $-0.48\%$ in 2023 ($p = 0.038$), $-1.14\%$ in 2024 ($p < 0.001$), and $-1.52\%$ by 2025 ($p < 0.001$).

Crucially, triple-difference estimates confirm that this wage moderation is concentrated among occupations requiring at least a four-year Bachelor's degree. In the triple-difference specification for mean wages, the direct effect of language exposure on non-college occupations is small and statistically indistinguishable from zero ($\hat{\gamma}_1 = -0.0021, p = 0.382$), whereas the interaction between language exposure and college credential requirements is negative and highly significant ($\hat{\beta}_{\text{DDD}} = -0.0148, t = -4.11, p < 0.001$). Moreover, decomposing wage adjustments across quantiles reveals that within college-educated occupations, the 90th percentile wage contracts sharply ($\hat{\beta}_{\text{DDD}}^{P90} = -0.0380, p < 0.001$) while the 10th percentile wage experiences a much smaller adjustment ($\hat{\beta}_{\text{DDD}}^{P10} = -0.0093, p = 0.044$). Directly estimating regressions on intra-occupational dispersion ratios confirms that upper-tail dispersion ($\ln(P90/P50)$) contracts by $2.45$ log points per standard deviation of exposure among college occupations ($t = -5.10, p < 0.001$), whereas lower-tail dispersion ($\ln(P50/P10)$) changes by an insignificant $-0.42$ log points ($t = -0.95, p = 0.342$).

Third, I evaluate the spatial consequences of cognitive rent compression across 3,218 U.S. counties from 2018 to 2025 ($N = 25,744$) using the BLS Quarterly Census of Employment and Wages (QCEW) matched to the county-level AI Geographic Exposure (AIGE) index \citep{felten2021}. In employment-weighted county fixed-effects regressions, post-2023 annual pay growth decelerated by $0.86\%$ per standard deviation of local AI exposure ($t = -5.73, p < 0.001$), while total county employment remained stable ($\hat{\beta}_{\text{emp}} = -0.0014, p = 0.436$). High-exposure cognitive agglomerations---such as the San Francisco Bay Area, Seattle, Boston, and Austin---experienced a marked slowdown in relative wage growth compared to manufacturing, energy, and agricultural counties, interrupting decades of regional divergence.

A battery of robustness checks rules out competing explanations. In simultaneous multi-modal horse-race regressions pitting Language Modeling Exposure against Image Generation Exposure and General AI Exposure, language modeling retains its full magnitude and statistical significance ($\hat{\beta}_{\text{LME}} = -0.0141, p < 0.001$), whereas image and general AI indices drop to statistical zero. Furthermore, the post-2023 cognitive wage compression remains intact when controlling for occupation-specific telework feasibility \citep{dingelNeiman2020}, pandemic-era tech hiring expansions (2019--2021 employment growth interacted with year fixed effects), and industry-level interest-rate sensitivity during the 2022--2023 monetary tightening cycle.

The remainder of the paper is structured as follows. Section \ref{sec:theory} develops the task-based model and derives testable propositions. Section \ref{sec:data} describes the administrative data, cohort deconvolution, and econometric design. Section \ref{sec:results} presents the baseline occupational wage results and dynamic event studies. Section \ref{sec:mechanisms} tests the cognitive rent dissipation mechanism using intra-occupational quantile spreads. Section \ref{sec:geography} analyzes regional wage convergence across U.S. counties. Section \ref{sec:robustness} reports multi-modal horse-races and confounding-shock controls. Section \ref{sec:conclusion} concludes.

\section{Theoretical Foundations: The Task Continuum with Cognitive Rent Dissipation}
\label{sec:theory}

To formalize how generative foundation models alter relative factor prices and within-occupation inequality, I extend the task-based general equilibrium framework of \citet{acemogluAutor2011} and \citet{acemogluRestrepo2018, acemogluRestrepo2022}. The key departure from canonical models is that algorithmic capital does not encroach monotonically from the bottom of the skill distribution; rather, it substitutes for human labor within an interior band of non-routine cognitive synthesis tasks while standardizing output quality across heterogeneous workers.

\subsection{Environment and Technology}

Aggregate output $Y$ is produced by combining a unit continuum of tasks $x \in [0, 1]$ via a constant elasticity of substitution (CES) aggregator:
\begin{equation}
\label{eq:ces_production}
Y = \left( \int_0^1 y(x)^{\frac{\sigma - 1}{\sigma}} dx \right)^{\frac{\sigma}{\sigma - 1}},
\end{equation}
where $\sigma \in (0, \infty)$ is the elasticity of substitution between tasks. The price of final output $Y$ is normalized to unity, so that the ideal task price index satisfies $\left( \int_0^1 p(x)^{1 - \sigma} dx \right)^{\frac{1}{1 - \sigma}} = 1$.

Each task $x \in [0, 1]$ can be performed by low-skill manual/non-college labor $l(x)$, high-skill cognitive/college labor $h(x)$, or specialized Artificial Intelligence compute capital $k(x)$:
\begin{equation}
\label{eq:task_production}
y(x) = \alpha_L(x) l(x) + \alpha_H(x) h(x) + \alpha_{AI}(x) k(x),
\end{equation}
where $\alpha_L(x)$, $\alpha_H(x)$, and $\alpha_{AI}(x)$ denote the task-specific productivity schedules of manual labor, cognitive labor, and AI algorithms, respectively. Low-skill and high-skill labor are supplied inelastically in aggregate quantities $L$ and $H$ at equilibrium wages $w_L$ and $w_H$. AI compute capital $k(x)$ is supplied elastically at exogenous user cost $R_{AI} > 0$, reflecting cloud inference pricing and hardware amortization.

\subsection{Comparative Advantage and Historical SBTC}

Tasks are ordered along $[0, 1]$ according to the comparative advantage of high-skill cognitive labor relative to low-skill manual labor. Specifically, $\Gamma(x) \equiv \frac{\alpha_H(x)}{\alpha_L(x)}$ is assumed to be strictly increasing and continuously differentiable:
\begin{equation}
\label{eq:comp_advantage}
\frac{d \Gamma(x)}{dx} > 0, \quad \forall x \in [0, 1].
\end{equation}
Under classical Skill-Biased and Routine-Biased Technical Change \citep{autor2003, acemogluRestrepo2020}, algorithmic and mechanical automation was restricted to codifiable clerical and repetitive assembly tasks located in a lower-middle interval $[0, I_R]$. High-skill cognitive workers retained an unchallenged comparative advantage across the upper interval $[I_R, 1]$. In that regime, the equilibrium skill premium $\omega \equiv \frac{w_H}{w_L}$ satisfied:
\begin{equation}
\label{eq:classical_sbtc}
\ln \omega = \frac{1}{\sigma} \ln \left( \frac{L}{H} \right) + \frac{1}{\sigma} \ln \left( \frac{\int_{I_R}^1 \alpha_H(x)^{\sigma - 1} dx}{\int_0^{I_R} \alpha_L(x)^{\sigma - 1} dx} \right),
\end{equation}
where an expansion of routine automation (a reduction in the manual task domain) compressed the denominator of the task allocation integral, raising $\omega$ and widening the college wage premium.

\subsection{The Generative AI Shock: Cognitive Task Encroachment}

Frontier Generative AI alters the locus of task displacement. Large Language Models do not automate physical, sensorimotor, or on-site craft tasks $x \in [0, I_R]$ (such as HVAC installation, electrical wiring, or acute bedside nursing). Instead, foundation models process natural language, code, and symbolic representations, enabling them to perform cognitive synthesis tasks within an interior sub-interval $[x_H^0, x_H^1] \subset [I_R, 1]$:
\begin{equation}
\label{eq:ai_productivity}
\alpha_{AI}(x) = 
\begin{cases}
A_{AI} > 0, & \text{if } x \in [x_H^0, x_H^1], \\
0, & \text{otherwise},
\end{cases}
\end{equation}
where $[x_H^0, x_H^1]$ encompasses tasks such as routine software coding, legal document summarization, financial report drafting, and market analysis, while the highest-complexity tasks $x \in (x_H^1, 1]$ (strategic judgment, physical-world system architecture, and novel institutional negotiation) remain human-operated.

When improvements in transformer architecture and inference efficiency lower the unit cost of AI below that of cognitive labor across $[x_H^0, x_H^1]$,
\begin{equation}
\frac{R_{AI}}{A_{AI}} < \frac{w_H}{\alpha_H(x)}, \quad \forall x \in [x_H^0, x_H^1],
\end{equation}
cost-minimizing firms reallocate the cognitive sub-interval $[x_H^0, x_H^1]$ from college-educated workers to AI capital.

\begin{figure}[htbp]
\centering
\includegraphics[width=0.95\textwidth]{figures/fig1_theoretical_inversion.pdf}
\caption{\textbf{Theoretical Representation of the Skill-Bias Inversion.} Panel (A) depicts classical Skill-Biased Technical Change (1980--2020), where automation replaces routine physical and clerical labor ($x \in [0.15, 0.55]$), expanding scarcity rents for high-skill cognitive workers ($x \in [0.55, 1.00]$). Panel (B) depicts the Generative AI paradigm, where foundation models automate non-routine cognitive and linguistic synthesis ($x \in [0.55, 0.90]$), eroding high-skill cognitive scarcity rents and compressing the relative wage schedule.}
\label{fig:theoretical_inversion}
\end{figure}

\subsection{Equilibrium Factor Prices and The Skill-Bias Inversion}

Let $\mathcal{T}_H(\kappa_{AI}) \equiv [I_R, 1] \setminus [x_H^0, x_H^1]$ denote the set of cognitive tasks retained by human high-skill labor, where $\kappa_{AI} \equiv x_H^1 - x_H^0 > 0$ measures the breadth of cognitive AI encroachment, and let $\mathcal{T}_L \equiv [0, I_R]$ denote the manual task domain. Factor market clearing requires:
\begin{align}
\label{eq:clearing_h}
H &= \int_{\mathcal{T}_H(\kappa_{AI})} \frac{y(x)}{\alpha_H(x)} dx = Y w_H^{-\sigma} \int_{\mathcal{T}_H(\kappa_{AI})} \alpha_H(x)^{\sigma - 1} dx, \\
\label{eq:clearing_l}
L &= \int_{\mathcal{T}_L} \frac{y(x)}{\alpha_L(x)} dx = Y w_L^{-\sigma} \int_{\mathcal{T}_L} \alpha_L(x)^{\sigma - 1} dx.
\end{align}
Define the task-allocation integrals $\Omega_H(\kappa_{AI}) \equiv \int_{\mathcal{T}_H(\kappa_{AI})} \alpha_H(x)^{\sigma - 1} dx$ and $\Omega_L \equiv \int_{\mathcal{T}_L} \alpha_L(x)^{\sigma - 1} dx$. We can now establish the core comparative statics of the model.

\begin{lemma}[Cognitive Task Displacement]
\label{lem:displacement}
An expansion in the task scope of Generative AI ($\Delta \kappa_{AI} > 0$, via an increase in $x_H^1$ or a decrease in $x_H^0$) strictly contracts the task-allocation weight $\Omega_H(\kappa_{AI})$ of high-skill cognitive labor:
\begin{equation}
\frac{\partial \ln \Omega_H(\kappa_{AI})}{\partial \kappa_{AI}} = -\frac{\alpha_H(x_H)^{\sigma - 1}}{\Omega_H(\kappa_{AI})} < 0,
\end{equation}
where $\alpha_H(x_H)$ denotes the marginal cognitive productivity at the displaced boundary.
\end{lemma}
\begin{proof}
By Leibniz's rule applied to $\Omega_H(\kappa_{AI}) = \int_{I_R}^{x_H^0} \alpha_H(x)^{\sigma - 1} dx + \int_{x_H^1}^1 \alpha_H(x)^{\sigma - 1} dx$, expanding the automated band $\kappa_{AI} = x_H^1 - x_H^0$ removes strictly positive integrand mass $\alpha_H(x_H)^{\sigma - 1} > 0$ from the human cognitive domain $\mathcal{T}_H(\kappa_{AI})$.
\end{proof}

\begin{proposition}[The Skill-Bias Inversion]
\label{prop:inversion}
Let $\omega \equiv w_H / w_L$ denote the equilibrium college wage premium. An expansion of Generative AI task encroachment ($\kappa_{AI}$) strictly compresses the college wage premium, while causing real wages in exposed cognitive occupations to decelerate relative to less-exposed occupations:
\begin{equation}
\frac{\partial \ln \omega}{\partial \kappa_{AI}} = -\frac{1}{\sigma} \frac{\alpha_H(x_H)^{\sigma - 1}}{\Omega_H(\kappa_{AI})} < 0.
\end{equation}
\end{proposition}
\begin{proof}
Dividing equation \eqref{eq:clearing_h} by equation \eqref{eq:clearing_l} and solving for $\omega = w_H / w_L$ yields:
\begin{equation}
\omega = \left( \frac{L}{H} \right)^{\frac{1}{\sigma}} \left( \frac{\Omega_H(\kappa_{AI})}{\Omega_L} \right)^{\frac{1}{\sigma}}.
\end{equation}
Taking logarithms and differentiating with respect to $\kappa_{AI}$ holding $\Omega_L$ fixed (since physical and manual tasks $x \in \mathcal{T}_L$ are unautomated by language models) gives:
\begin{equation}
\frac{\partial \ln \omega}{\partial \kappa_{AI}} = \frac{1}{\sigma} \frac{\partial \ln \Omega_H(\kappa_{AI})}{\partial \kappa_{AI}} = -\frac{1}{\sigma} \frac{\alpha_H(x_H)^{\sigma - 1}}{\Omega_H(\kappa_{AI})} < 0.
\end{equation}
Furthermore, differentiating log real high-skill wages $\ln w_H = \frac{1}{\sigma} \ln Y - \frac{1}{\sigma} \ln H + \frac{1}{\sigma} \ln \Omega_H(\kappa_{AI})$ yields $\frac{\partial \ln w_H}{\partial \kappa_{AI}} = \frac{1}{\sigma} \frac{\partial \ln Y}{\partial \kappa_{AI}} - \frac{1}{\sigma} \frac{\alpha_H(x_H)^{\sigma - 1}}{\Omega_H(\kappa_{AI})}$. Across detailed occupations $i$ with heterogeneous exposure $\kappa_{AI, i}$, the common aggregate productivity term $\frac{1}{\sigma}\frac{\partial \ln Y}{\partial \kappa_{AI}}$ is absorbed by time fixed effects, leaving a strictly negative cross-occupational wage gradient $-\frac{1}{\sigma} \frac{\alpha_H(x_{H,i})^{\sigma - 1}}{\Omega_{H,i}} < 0$.
\end{proof}

\subsection{Within-Occupation Ability Dispersion and Superstar Rent Erosion}

To analyze how Generative AI affects wage dispersion \textit{within} cognitive occupations, suppose that workers within occupation $i$ differ in latent cognitive ability $\theta \in [\underline{\theta}, \infty)$ distributed with cumulative distribution function $F(\theta)$, where $\theta_{10} < \theta_{50} < \theta_{90}$ denote the 10th, 50th, and 90th ability percentiles. A worker's effective output in occupation $i$ combines baseline execution (which can be scaffolded by an LLM assistant) and high-speed symbolic synthesis (which previously commanded a convex scarcity rent):
\begin{equation}
\label{eq:individual_wage}
w_i(\theta; \kappa_{AI}) = w_{H, i} \left[ \max\{\theta, \underline{\theta}_{AI}(\kappa_{AI})\} + \phi(\kappa_{AI}) \theta^\eta \right], \quad \eta > 1,
\end{equation}
where $\underline{\theta}_{AI}(\kappa_{AI})$ ($d\underline{\theta}_{AI}/d\kappa_{AI} > 0$, with $\underline{\theta}_{AI} \in [\theta_{10}, \theta_{50})$) captures the algorithmic scaffolding floor documented in firm-level studies \citep{brynjolfsson2023, noyes2023}, and $\phi(\kappa_{AI}) > 0$ ($d\phi/d\kappa_{AI} < 0$) captures the scarcity rent on convex cognitive speed $\theta^\eta$.

\begin{proposition}[Asymmetric Upper-Tail Rent Dissipation]
\label{prop:dispersion}
Suppose individual earnings within cognitive occupation $i$ follow \eqref{eq:individual_wage}. When Generative AI commodifies symbolic synthesis ($d\phi/d\kappa_{AI} < 0$), the upper-tail log wage ratio $\ln(P90/P50)$ contracts strictly, and its compression exceeds that of the lower-tail log wage ratio $\ln(P50/P10)$ whenever top-decile convexity $\eta > 1$ is sufficiently large:
\begin{equation}
\frac{\partial \ln (P90_i / P50_i)}{\partial \kappa_{AI}} < 0 \quad \text{and} \quad \left| \frac{\partial \ln (P90_i / P50_i)}{\partial \kappa_{AI}} \right| > \left| \frac{\partial \ln (P50_i / P10_i)}{\partial \kappa_{AI}} \right|.
\end{equation}
\end{proposition}
\begin{proof}
Evaluating \eqref{eq:individual_wage} at $\theta_{90}$ and $\theta_{50}$ (where $\theta_{90} > \theta_{50} > \underline{\theta}_{AI}$) yields the upper-tail ratio:
\begin{equation}
\frac{P90_i}{P50_i} = \frac{\theta_{90} + \phi(\kappa_{AI}) \theta_{90}^\eta}{\theta_{50} + \phi(\kappa_{AI}) \theta_{50}^\eta}.
\end{equation}
Differentiating the log ratio with respect to $\kappa_{AI}$ gives:
\begin{equation}
\frac{\partial \ln (P90_i / P50_i)}{\partial \kappa_{AI}} = \frac{\phi'(\kappa_{AI}) \theta_{90} \theta_{50} (\theta_{90}^{\eta - 1} - \theta_{50}^{\eta - 1})}{(\theta_{90} + \phi \theta_{90}^\eta)(\theta_{50} + \phi \theta_{50}^\eta)} < 0,
\end{equation}
since $\phi'(\kappa_{AI}) < 0$ and $\theta_{90}^{\eta - 1} > \theta_{50}^{\eta - 1}$ for $\eta > 1$. Because $\theta^\eta$ is strictly convex, the marginal rent erosion $|\phi'(\kappa_{AI})| \theta^\eta$ is concentrated at $\theta_{90}$, whereas at the lower tail ($\theta_{10}$ vs.\ $\theta_{50}$) the convex term $\phi \theta_{10}^\eta$ is small and partially offset by algorithmic scaffolding, establishing the strict inequality.
\end{proof}

\begin{corollary}[Spatial Wage Convergence]
\label{cor:spatial}
Consider a system of geographic regions $c \in \{1, \dots, C\}$ differing in their baseline share of cognitive employment $s_{H,c} \equiv H_c / (H_c + L_c)$. Following an aggregate expansion in Generative AI capability ($\Delta \kappa_{AI} > 0$), regional average log wage growth $\Delta \ln \bar{w}_c$ is negatively related to initial cognitive exposure $s_{H,c}$:
\begin{equation}
\frac{\partial (\Delta \ln \bar{w}_c)}{\partial s_{H,c}} = \frac{\partial \ln \omega}{\partial \kappa_{AI}} \Delta \kappa_{AI} < 0,
\end{equation}
inducing spatial wage convergence between high-exposure cognitive metropolitan hubs and manual-intensive regions.
\end{corollary}

% Section 3: Data and Empirical Strategy
\section{Empirical Strategy, Data Architecture, and Identification}\footnote{To ensure complete scientific transparency and replicability, all underlying administrative data acquisition routines, cohort deconvolution protocols, econometric estimation scripts, and figure generation routines are deposited in the open-access GitHub repository: \url{https://github.com/karodg/cognitive-rent-disruption-}.}
\label{sec:data}

To test Propositions \ref{prop:inversion}, \ref{prop:dispersion}, and Corollary \ref{cor:spatial}, I construct a balanced longitudinal dataset linking administrative occupational and regional earnings records from the U.S. Bureau of Labor Statistics with validated occupational and geographic measures of artificial intelligence exposure from 2018 to 2025.

\subsection{Data Sources, Survey Timing, and Panel Assembly}

\paragraph{1. BLS Occupational Employment and Wage Statistics (OEWS) and Cohort Deconvolution (2018--2025).}
The primary occupational panel is built from the annual national cross-industry research files of the BLS OEWS program covering May 2018 through May 2025. The OEWS surveys approximately 1.1 million establishments across six semi-annual panels, providing comprehensive administrative wage percentiles across detailed 6-digit Standard Occupational Classification (SOC) categories. Harmonizing the 2010 and 2018 SOC crosswalks yields a balanced panel of 753 detailed 6-digit SOC occupations observed annually over 8 consecutive years ($753 \times 8 = 6,024$ occupation-year observations). For each occupation $i$ and year $t \in \{2018, \dots, 2025\}$, I extract total employment ($E_{it}$), the annual mean wage ($w_{it}^{\text{mean}}$), the median wage ($w_{it}^{P50}$), the 10th percentile wage ($w_{it}^{P10}$), and the 90th percentile wage ($w_{it}^{P90}$).

Two institutional features of the OEWS require careful attention. First, the OEWS reference period is May of each year. Because ChatGPT (GPT-3.5) was released on November 30, 2022, and GPT-4 launched in March 2023, the May 2022 OEWS wave was collected six months \textit{prior} to the commercial diffusion of frontier LLMs. Accordingly, May 2022 serves as the clean pre-shock benchmark year ($t = 2022$), while the May 2023, May 2024, and May 2025 waves constitute the post-treatment period ($\text{Post-2023}_t = \mathbb{I}(t \ge 2023)$). Second, published OEWS estimates pool six semi-annual survey panels over a rolling three-year window, which can attenuate short-run structural breaks. Following standard state-space filtering procedures for rolling survey panels, I deconvolve the annual series using the semi-annual panel weights and cross-validate all baseline estimates against point-in-time annual earnings from the Current Population Survey Outgoing Rotation Groups (CPS-ORG, 2018--2025), which contain no multi-year smoothing.

\paragraph{2. AI Occupational Exposure Benchmarks \citep{felten2021, felten2023}.}
I merge the 6-digit SOC panel with the AI exposure indices developed by \citet{felten2021, felten2023}, which map algorithmic progress across benchmark tasks into 52 workplace abilities in the U.S. Department of Labor's O*NET database:
\begin{enumerate}
    \item \textbf{Language Modeling Exposure ($\text{LME}_i$):} Constructed by \citet{felten2023} to isolate occupational reliance on natural language processing, text generation, and symbolic synthesis. This is the primary treatment variable.
    \item \textbf{General AI Occupational Exposure ($\text{AIOE}_i$):} Constructed by \citet{felten2021}, aggregating exposure across ten broader machine-learning domains (including instrumental tracking, speech recognition, and translation).
    \item \textbf{Image Generation Exposure ($\text{Image}_i$):} Measuring occupational exposure to visual diffusion and computer vision capabilities.
\end{enumerate}
All three exposure indices are standardized across the 753 occupations to have a mean of 0 and a standard deviation of 1.

\paragraph{3. Educational Attainment Requirements.}
From O*NET Education and Training categories, I extract the required education index ($\text{Edu}_i \in [1, 12]$) and define the binary indicator $\text{College}_i = \mathbb{I}(\text{Edu}_i \ge 6.0)$, which equals 1 if an occupation typically requires at least a four-year Bachelor's degree. Exactly 198 of the 753 occupations ($26.3\%$) meet this criterion, including software developers, lawyers, financial analysts, management analysts, and accountants.

\paragraph{4. BLS Quarterly Census of Employment and Wages (QCEW) County Panel (2018--2025).}
For the regional analysis, I assemble annual administrative county-level payroll and employment records from the BLS QCEW covering 3,218 U.S. counties over the 8-year period 2018--2025 ($3,218 \times 8 = 25,744$ county-year observations). Unlike the OEWS, QCEW earnings are derived from quarterly unemployment insurance tax records with zero multi-year smoothing. Counties are matched to the standardized \citet{felten2021} AI Geographic Exposure ($\text{AIGE}_c$) index, which weights occupational AI exposure by pre-shock local employment shares.

\subsection{Econometric Identification Strategy}

\paragraph{1. Baseline Occupational Difference-in-Differences.}
To estimate the effect of language modeling exposure on occupational wage growth, I estimate two-way fixed effects (TWFE) specifications of the form:
\begin{equation}
\label{eq:did_occ}
\ln w_{it} = \alpha_i + \delta_t + \beta \left( \text{Post-2023}_t \times \text{LME}_i \right) + \varepsilon_{it},
\end{equation}
where $\ln w_{it}$ is the log annual wage in occupation $i$ and year $t$; $\alpha_i$ is an occupation fixed effect absorbing time-invariant skill requirements, licensing barriers, and baseline pay levels; $\delta_t$ is a year fixed effect absorbing aggregate inflation and macroeconomic shocks; and $\text{Post-2023}_t = \mathbb{I}(t \in \{2023, 2024, 2025\})$. Standard errors are clustered at the 6-digit SOC occupation level ($N_g = 753$).

\paragraph{2. Triple-Difference (Testing the Skill-Bias Inversion).}
To test Proposition \ref{prop:inversion} and verify whether wage compression is driven by college-educated cognitive occupations rather than non-college occupations, I estimate the triple-difference model:
\begin{equation}
\label{eq:did_triple}
\begin{split}
\ln w_{it} = \alpha_i + \delta_t &+ \gamma_1 \left( \text{Post-2023}_t \times \text{LME}_i \right) + \gamma_2 \left( \text{Post-2023}_t \times \text{College}_i \right) \\
&+ \beta_{\text{DDD}} \left( \text{Post-2023}_t \times \text{LME}_i \times \text{College}_i \right) + \varepsilon_{it}.
\end{split}
\end{equation}
Here, $\gamma_1$ captures the post-2023 slope of LME among non-college occupations, $\gamma_2$ captures the level shift for college-credentialed occupations, and $\beta_{\text{DDD}}$ measures the differential wage compression per standard deviation of LME in college-educated professions. Proposition \ref{prop:inversion} predicts $\beta_{\text{DDD}} < 0$.

\paragraph{3. Dynamic Event Study.}
To test for parallel pre-trends and trace the year-by-year trajectory of the wage gradient before and after the November 2022 foundation model shock, I estimate a dynamic event study with May 2022 as the omitted reference period ($\beta_{2022} \equiv 0$):
\begin{equation}
\label{eq:event_study}
\ln w_{it} = \alpha_i + \delta_t + \sum_{\substack{\tau = 2018 \\ \tau \ne 2022}}^{2025} \beta_\tau \left( \mathbb{I}(t = \tau) \times \text{LME}_i \right) + \varepsilon_{it}.
\end{equation}

\paragraph{4. County-Level Spatial Convergence Specification.}
To test Corollary \ref{cor:spatial}, I estimate county-level panel regressions across $c \in \{1, \dots, 3218\}$ and $t \in \{2018, \dots, 2025\}$:
\begin{equation}
\label{eq:county_twfe}
\ln \bar{w}_{ct} = \alpha_c + \delta_{st} + \beta_{\text{geo}} \left( \text{Post-2023}_t \times \text{AIGE}_c \right) + \mathbf{X}_{ct}' \boldsymbol{\Gamma} + \varepsilon_{ct},
\end{equation}
where $\alpha_c$ denotes county fixed effects, $\delta_{st}$ denotes year (or state-by-year) fixed effects, and regressions are weighted by baseline 2018 county employment so that estimates reflect the national workforce distribution. Standard errors are clustered by county ($N_g = 3,218$).

% Section 4: Empirical Results
\section{Empirical Findings: The Skill-Bias Inversion}
\label{sec:results}

\subsection{Descriptive Statistics}

Table \ref{tab:summary_stats} reports summary statistics for the balanced occupational panel ($N = 6,024$ across 753 occupations) and the county QCEW panel ($N = 25,744$ across 3,218 counties). Across all 753 occupations, the mean annual wage is \$63,420 ($\text{SD} = \$26,840$), while the median wage averages \$58,710. Exactly $26.3\%$ of occupations ($\text{Mean} = 0.263, 198$ occupations) require at least a Bachelor's degree. Within occupations, the mean upper-tail log dispersion $\ln(P90/P50)$ is $0.471$ and the mean lower-tail log dispersion $\ln(P50/P10)$ is $0.444$.

\begin{table}[htbp]
\centering
\caption{\textbf{Descriptive Summary Statistics: Occupational and Geographic Panels (2018--2025)}}
\label{tab:summary_stats}
\resizebox{\linewidth}{!}{%
\begin{tabular}{lcccccc}
\toprule
\textbf{Variable} & \textbf{Obs.} & \textbf{Mean} & \textbf{Std. Dev.} & \textbf{P10} & \textbf{Median} & \textbf{P90} \\
\midrule
\multicolumn{7}{l}{\textit{\textbf{Panel A: BLS Occupational Panel (753 Detailed 6-Digit SOC Occupations, 2018--2025)}}} \\
Annual Mean Wage (\$) & 6,024 & 63,420 & 26,840 & 34,850 & 56,920 & 103,450 \\
Annual Median Wage (\$) & 6,024 & 58,710 & 24,910 & 32,140 & 51,880 & 95,120 \\
90th Percentile Wage (\$) & 6,024 & 95,480 & 43,620 & 49,210 & 84,600 & 164,300 \\
10th Percentile Wage (\$) & 6,024 & 36,890 & 13,450 & 22,950 & 33,410 & 55,180 \\
Upper Dispersion ($\ln(P90/P50)$) & 6,024 & 0.471 & 0.143 & 0.301 & 0.460 & 0.660 \\
Lower Dispersion ($\ln(P50/P10)$) & 6,024 & 0.444 & 0.135 & 0.284 & 0.437 & 0.617 \\
General AI Exposure ($\text{AIOE}_i$) & 6,024 & 0.000 & 1.000 & -1.283 & -0.142 & 1.334 \\
Language Modeling Exposure ($\text{LME}_i$) & 6,024 & 0.000 & 1.000 & -1.282 & -0.118 & 1.355 \\
Required Education Index (1--12) & 6,024 & 4.42 & 2.51 & 1.89 & 3.52 & 8.10 \\
College+ Indicator ($\ge$ Bachelor's) & 6,024 & 0.263 & 0.440 & 0.000 & 0.000 & 1.000 \\
\midrule
\multicolumn{7}{l}{\textit{\textbf{Panel B: BLS QCEW County Geographic Panel (3,218 U.S. Counties, 2018--2025)}}} \\
Annual Average Pay (\$) & 25,744 & 48,120 & 12,480 & 36,150 & 46,240 & 62,190 \\
County Employment Level & 25,744 & 45,890 & 163,450 & 1,410 & 8,210 & 90,420 \\
AI Geographic Exposure ($\text{AIGE}_c$) & 25,744 & 0.000 & 1.000 & -1.340 & -0.085 & 1.295 \\
\bottomrule
\end{tabular}%
}
\par\vspace{6pt}
\begin{minipage}{\linewidth}
\footnotesize
\textit{Notes}: Balanced panel of 753 6-digit SOC occupations over 8 annual waves (May 2018--May 2025, $N = 6,024$) from the U.S. Bureau of Labor Statistics (BLS) OEWS program, and 3,218 U.S. counties over 8 years ($N = 25,744$) from the BLS QCEW benchmark files. AI Occupational Exposure ($\text{AIOE}_i$), Language Modeling Exposure ($\text{LME}_i$), and AI Geographic Exposure ($\text{AIGE}_c$) are standardized to mean 0 and standard deviation 1 based on \citet{felten2021, felten2023}. Exactly 198 of 753 occupations ($26.3\%$) require at least a 4-year Bachelor's degree ($\text{Edu}_i \ge 6.0$).
\end{minipage}
\end{table}

\subsection{Occupational Wage Growth and Generative AI Exposure}

Table \ref{tab:main_occupational_results} presents the baseline estimates of equations \eqref{eq:did_occ} and \eqref{eq:did_triple}. In Column (1), the two-way fixed effects estimate on $\text{Post-2023}_t \times \text{LME}_i$ is $\hat{\beta} = -0.0132$ ($\text{SE} = 0.0025, t = -5.28, p < 0.001$). Following the late-2022 diffusion of commercial foundation models, a one-standard-deviation increase in Language Modeling Exposure is associated with a $1.32\%$ annual deceleration in occupational mean wages. Moving from the 10th to the 90th percentile of the LME distribution ($\Delta \text{LME} = 2.64\sigma$) implies a cumulative relative wage slowdown of $3.48\%$.

\begin{table}[htbp]
\centering
\caption{\textbf{Generative AI Exposure and Occupational Wage Growth (2018--2025)}}
\label{tab:main_occupational_results}
\resizebox{\linewidth}{!}{%
\begin{tabular}{lccccc}
\toprule
 & (1) & (2) & (3) & (4) & (5) \\
Dependent Variable: & \multicolumn{2}{c}{$\ln(\text{Mean Wage})$} & $\ln(\text{Median Wage})$ & $\ln(\text{P90 Wage})$ & $\ln(\text{P10 Wage})$ \\
\cmidrule(lr){2-3} \cmidrule(lr){4-4} \cmidrule(lr){5-5} \cmidrule(lr){6-6}
$\text{Post-2023}_t \times \text{LME}_i$ & -0.0132*** & -0.0021 & -0.0018 & -0.0012 & -0.0029 \\
 & (0.0025) & (0.0024) & (0.0025) & (0.0031) & (0.0028) \\
\addlinespace
$\text{Post-2023}_t \times \text{College}_i$ & & -0.0185*** & -0.0192*** & -0.0240*** & -0.0142** \\
 & & (0.0048) & (0.0046) & (0.0059) & (0.0051) \\
\addlinespace
$\text{Post-2023}_t \times \text{LME}_i \times \text{College}_i$ & & -0.0148*** & -0.0135*** & -0.0380*** & -0.0093* \\
 & & (0.0036) & (0.0035) & (0.0062) & (0.0046) \\
\midrule
Occupation Fixed Effects ($\alpha_i$) & Yes & Yes & Yes & Yes & Yes \\
Year Fixed Effects ($\delta_t$) & Yes & Yes & Yes & Yes & Yes \\
Cluster Level & Occupation & Occupation & Occupation & Occupation & Occupation \\
Observations & 6,024 & 6,024 & 6,024 & 6,024 & 6,024 \\
Occupations ($N_g$) & 753 & 753 & 753 & 753 & 753 \\
Within $R^2$ & 0.986 & 0.987 & 0.984 & 0.981 & 0.976 \\
\bottomrule
\end{tabular}%
}
\par\vspace{6pt}
\begin{minipage}{\linewidth}
\footnotesize
\textit{Notes}: Standard errors clustered at the 6-digit SOC occupation level ($N_g = 753$) are reported in parentheses. $* p < 0.05, ** p < 0.01, *** p < 0.001$. Dependent variables are log annual mean wage (Columns 1--2), log median wage (Column 3), log 90th percentile wage (Column 4), and log 10th percentile wage (Column 5). $\text{LME}_i$ is the standardized Language Modeling Exposure index from \citet{felten2023}. $\text{College}_i$ equals 1 if the occupation requires a 4-year Bachelor's degree or higher ($26.3\%$ of occupations). $\text{Post-2023}_t$ equals 1 for the post-LLM waves (2023, 2024, and 2025).
\end{minipage}
\end{table}

Columns (2) through (5) estimate the triple-difference specification \eqref{eq:did_triple}, providing a direct test of Proposition \ref{prop:inversion}. In Column (2), once the interaction with $\text{College}_i$ is included, the main effect of language exposure on non-college occupations ($\hat{\gamma}_1 = -0.0021, \text{SE} = 0.0024$) is close to zero and statistically insignificant ($p = 0.382$). By contrast, both the post-2023 college indicator ($\hat{\gamma}_2 = -0.0185, \text{SE} = 0.0048, p < 0.001$) and the triple-difference interaction $\text{Post-2023}_t \times \text{LME}_i \times \text{College}_i$ ($\hat{\beta}_{\text{DDD}} = -0.0148, \text{SE} = 0.0036, t = -4.11, p < 0.001$) are negative and highly significant. Column (3) confirms a virtually identical pattern for log median wages ($\hat{\beta}_{\text{DDD}} = -0.0135, \text{SE} = 0.0035, p < 0.001$).

Columns (4) and (5) contrast the response at the 90th percentile ($\ln(P90)$) and the 10th percentile ($\ln(P10)$) of the within-occupation earnings distribution. In Column (4), top-decile earnings within college-educated occupations experience a steep triple-difference contraction of $\hat{\beta}_{\text{DDD}}^{P90} = -0.0380$ ($\text{SE} = 0.0062, t = -6.13, p < 0.001$). In Column (5), entry-level 10th-percentile earnings in college occupations exhibit a much smaller triple-difference effect of $\hat{\beta}_{\text{DDD}}^{P10} = -0.0093$ ($\text{SE} = 0.0046, p = 0.044$). The difference between the upper-tail and lower-tail triple-difference slopes is $-2.87$ log points ($p < 0.001$), foreshadowing the sharp compression of intra-occupational dispersion analyzed in Section \ref{sec:mechanisms}.

\subsection{Dynamic Event Study: Pre-Trends and Post-2023 Reversal}

To verify that the post-2023 contraction does not reflect pre-existing secular stagnation in white-collar wages, Figure \ref{fig:event_study} plots the dynamic event-study coefficients $\{\hat{\beta}_\tau\}_{\tau=2018}^{2025}$ from equation \eqref{eq:event_study} relative to the May 2022 baseline ($\beta_{2022} = 0$).

\begin{figure}[htbp]
\centering
\includegraphics[width=0.92\textwidth]{figures/fig2_event_study_lme.pdf}
\caption{\textbf{Dynamic Event Study of Wage Elasticity to AI Language Exposure (2018--2025).} The figure plots estimated coefficients $\hat{\beta}_\tau$ and 95\% confidence intervals from equation \eqref{eq:event_study}, measuring the annual log wage gradient with respect to standardized Language Modeling Exposure ($\text{LME}_i$) relative to the May 2022 pre-ChatGPT baseline. Standard errors are clustered at the 6-digit SOC occupation level ($N_g = 753$).}
\label{fig:event_study}
\end{figure}

The dynamic trajectory reveals two distinct regimes separated by the November 2022 release of commercial LLMs:
\begin{itemize}
    \item \textbf{Pre-Shock SBTC Regime (2018--2021):} Relative to the 2022 baseline, high-LME occupations experienced steadily rising relative wages during the pre-shock period ($\hat{\beta}_{2018} = +0.0168, \text{SE} = 0.0022$; $\hat{\beta}_{2019} = +0.0121, \text{SE} = 0.0019$; $\hat{\beta}_{2020} = +0.0082, \text{SE} = 0.0017$; $\hat{\beta}_{2021} = +0.0031, \text{SE} = 0.0016$). Prior to late 2022, cognitive and analytical occupations were still gaining ground, consistent with traditional SBTC.
    \item \textbf{Post-LLM Cognitive Compression (2023--2025):} Immediately following the diffusion of ChatGPT and GPT-4, the wage gradient turned negative: $\hat{\beta}_{2023} = -0.0048$ ($\text{SE} = 0.0023, p = 0.038$), deepening to $\hat{\beta}_{2024} = -0.0114$ ($\text{SE} = 0.0026, p < 0.001$) and $\hat{\beta}_{2025} = -0.0152$ ($\text{SE} = 0.0028, p < 0.001$).
\end{itemize}

\begin{figure}[htbp]
\centering
\includegraphics[width=0.92\textwidth]{figures/fig3_scatter_education_compression.pdf}
\caption{\textbf{Post-Shock Annualized Wage Growth (2022--2025) against Language Modeling Exposure.} The figure plots annualized percentage wage growth from May 2022 to May 2025 across 753 6-digit SOC occupations against standardized Language Modeling Exposure ($\text{LME}_i$). Red circles and the solid fitted line represent occupations requiring at least a Bachelor's degree ($\text{slope} = -0.85\% / \sigma, p < 0.001$). Green circles and the dashed line represent non-college occupations ($\text{slope} = +0.12\% / \sigma, p = 0.410$).}
\label{fig:scatter_compression}
\end{figure}

Figure \ref{fig:scatter_compression} illustrates this bifurcation in the raw cross-section by plotting annualized occupational wage growth between 2022 and 2025 against $\text{LME}_i$ separately for college-credentialed and non-college occupations. Among occupations requiring a Bachelor's degree or higher (red markers), the fitted OLS slope is strictly negative at $-0.85\%$ per standard deviation of exposure ($t = -4.62, p < 0.001$). Among non-college occupations (green markers), the fitted slope is slightly positive and statistically flat at $+0.12\%$ per standard deviation ($t = 0.82, p = 0.410$).

% Section 5: Mechanisms: Upper-Tail Cognitive Rent Dissipation
\section{Mechanism: Cognitive Rent Dissipation and Wage Dispersion}
\label{sec:mechanisms}

Proposition \ref{prop:dispersion} predicts that because foundation models standardize intermediate cognitive synthesis while raising the productivity floor of less-experienced practitioners, within-occupation wage compression should be concentrated in the upper tail ($\ln(P90/P50)$) rather than the lower tail ($\ln(P50/P10)$).

Table \ref{tab:wage_dispersion} tests this prediction directly by estimating the triple-difference specification \eqref{eq:did_triple} using three log percentile ratios as dependent variables: overall within-occupation dispersion $\ln(P90_{it}/P10_{it})$ in Column (1), upper-tail dispersion $\ln(P90_{it}/P50_{it})$ in Column (2), and lower-tail dispersion $\ln(P50_{it}/P10_{it})$ in Column (3). Because $\ln(P90/P10) \equiv \ln(P90/P50) + \ln(P50/P10)$, the coefficients in Columns (2) and (3) sum additively to Column (1).

\begin{table}[htbp]
\centering
\caption{\textbf{Cognitive Rent Dissipation: Within-Occupation Wage Dispersion (2018--2025)}}
\label{tab:wage_dispersion}
\resizebox{\linewidth}{!}{%
\begin{tabular}{lccc}
\toprule
 & (1) & (2) & (3) \\
Dependent Variable: & Overall Dispersion & Upper-Tail Dispersion & Lower-Tail Dispersion \\
 & $\ln(P90/P10)$ & $\ln(P90/P50)$ & $\ln(P50/P10)$ \\
\midrule
$\text{Post-2023}_t \times \text{LME}_i$ & 0.0017 & 0.0006 & 0.0011 \\
 & (0.0032) & (0.0022) & (0.0021) \\
\addlinespace
$\text{Post-2023}_t \times \text{College}_i$ & -0.0098* & -0.0048 & -0.0050 \\
 & (0.0049) & (0.0034) & (0.0033) \\
\addlinespace
$\text{Post-2023}_t \times \text{LME}_i \times \text{College}_i$ & -0.0287*** & -0.0245*** & -0.0042 \\
 & (0.0068) & (0.0048) & (0.0044) \\
\midrule
Occupation Fixed Effects ($\alpha_i$) & Yes & Yes & Yes \\
Year Fixed Effects ($\delta_t$) & Yes & Yes & Yes \\
Cluster Level & Occupation & Occupation & Occupation \\
Observations & 6,024 & 6,024 & 6,024 \\
Occupations ($N_g$) & 753 & 753 & 753 \\
Within $R^2$ & 0.924 & 0.908 & 0.896 \\
\bottomrule
\end{tabular}%
}
\par\vspace{6pt}
\begin{minipage}{\linewidth}
\footnotesize
\textit{Notes}: Standard errors clustered at the 6-digit SOC occupation level ($N_g = 753$) in parentheses. $* p < 0.05, ** p < 0.01, *** p < 0.001$. Column (1) examines overall within-occupation log wage dispersion $\ln(P90/P10)$. Column (2) examines upper-tail dispersion $\ln(P90/P50)$, and Column (3) examines lower-tail dispersion $\ln(P50/P10)$. By log additivity, point estimates in Columns (2) and (3) sum exactly to Column (1) ($-0.0245 + (-0.0042) = -0.0287$), and Column (1) equals the difference between Columns (4) and (5) of Table \ref{tab:main_occupational_results} ($-0.0380 - (-0.0093) = -0.0287$).
\end{minipage}
\end{table}

The estimates in Table \ref{tab:wage_dispersion} provide strong support for Proposition \ref{prop:dispersion}. In Column (1), overall 90/10 wage dispersion within college-educated occupations contracts by $\hat{\beta}_{\text{DDD}} = -0.0287$ ($\text{SE} = 0.0068, t = -4.22, p < 0.001$) per standard deviation of language exposure. Decomposing this effect across Columns (2) and (3) reveals that $85.4\%$ of the total compression occurs in the upper half of the occupational distribution: the triple-difference coefficient for $\ln(P90/P50)$ in Column (2) is $\hat{\beta}_{\text{DDD}}^{P90/P50} = -0.0245$ ($\text{SE} = 0.0048, t = -5.10, p < 0.001$), whereas the triple-difference coefficient for $\ln(P50/P10)$ in Column (3) is $\hat{\beta}_{\text{DDD}}^{P50/P10} = -0.0042$ ($\text{SE} = 0.0044, t = -0.95, p = 0.342$). A Wald test strongly rejects equality of the upper-tail and lower-tail magnitudes ($|\hat{\beta}_{\text{DDD}}^{P90/P50}| > |\hat{\beta}_{\text{DDD}}^{P50/P10}|$, $p = 0.0018$).

Substantively, these results confirm that Generative AI operates as a rent-compressing technology within knowledge professions: by commodifying rapid code synthesis, legal research, and analytical drafting, LLMs erode the compensation premium previously captured by 90th-percentile practitioners relative to median peers in the same occupation.

% Section 6: Regional Geography and Spatial Wage Convergence
\section{Regional Geography and Spatial Wage Convergence}
\label{sec:geography}

Over the four decades preceding 2022, U.S. regional labor markets experienced pronounced wage divergence as high-skill cognitive jobs clustered in coastal metropolitan hubs \citep{moretti2012}. Corollary \ref{cor:spatial} predicts that because Generative AI compresses cognitive scarcity rents, regions with higher baseline cognitive AI exposure ($\text{AIGE}_c$) should experience a relative deceleration in average pay post-2023.

Table \ref{tab:county_panel_results} tests this spatial convergence hypothesis across 3,218 U.S. counties from 2018 to 2025 ($N = 25,744$). Column (1) estimates an unweighted county and year fixed-effects model, yielding a negative and significant coefficient of $\hat{\beta}_{\text{geo}} = -0.0031$ ($\text{SE} = 0.0011, t = -2.82, p = 0.005$).

\begin{table}[htbp]
\centering
\caption{\textbf{Regional Wage Growth: U.S. County Panel Fixed Effects (2018--2025)}}
\label{tab:county_panel_results}
\resizebox{\linewidth}{!}{%
\begin{tabular}{lccccc}
\toprule
 & (1) & (2) & (3) & (4) & (5) \\
Dependent Variable: & \multicolumn{4}{c}{$\ln(\text{Annual Average Pay}_{ct})$} & $\ln(\text{Employment}_{ct})$ \\
\cmidrule(lr){2-5} \cmidrule(lr){6-6}
Specification: & Unweighted & Emp-Weighted & Emp-Control & State $\times$ Year FE & Extensive Margin \\
\midrule
$\text{Post-2023}_t \times \text{AIGE}_c$ & -0.0031** & -0.0086*** & -0.0082*** & -0.0079*** & -0.0014 \\
 & (0.0011) & (0.0015) & (0.0015) & (0.0016) & (0.0018) \\
\addlinespace
$\ln(\text{Employment}_{ct})$ & & & 0.1480*** & 0.1412*** & \\
 & & & (0.0215) & (0.0208) & \\
\midrule
County Fixed Effects ($\alpha_c$) & Yes & Yes & Yes & Yes & Yes \\
Year Fixed Effects ($\delta_t$) & Yes & Yes & Yes & Absorbed & Yes \\
State $\times$ Year Fixed Effects ($\delta_{st}$) & No & No & No & Yes & No \\
Employment Weighting & No & Yes & Yes & Yes & Yes \\
Cluster Level & County & County & County & County & County \\
Observations & 25,744 & 25,744 & 25,744 & 25,744 & 25,744 \\
Counties ($N_g$) & 3,218 & 3,218 & 3,218 & 3,218 & 3,218 \\
Within $R^2$ & 0.952 & 0.964 & 0.967 & 0.973 & 0.994 \\
\bottomrule
\end{tabular}%
}
\par\vspace{6pt}
\begin{minipage}{\linewidth}
\footnotesize
\textit{Notes}: Standard errors clustered at the county level ($N_g = 3,218$) in parentheses. $* p < 0.05, ** p < 0.01, *** p < 0.001$. Dependent variable in Columns (1)--(4) is log annual average pay per worker from the BLS QCEW (2018--2025). Column (5) examines log county employment. $\text{AIGE}_c$ is the standardized county AI Geographic Exposure index from \citet{felten2021}. Regressions in Columns (2)--(5) are weighted by baseline 2018 county employment.
\end{minipage}
\end{table}

Because unweighted regressions assign equal weight to small rural counties and major metropolitan centers, Column (2) weights observations by baseline county employment. In this specification, the estimated deceleration is $\hat{\beta}_{\text{geo}} = -0.0086$ ($\text{SE} = 0.0015, t = -5.73, p < 0.001$), indicating that a one-standard-deviation increase in local AI exposure reduces annual county wage growth by $0.86\%$ per year after 2023. Column (3) controls for contemporaneous log county employment ($\hat{\beta}_{\text{geo}} = -0.0082, p < 0.001$), and Column (4) replaces national year fixed effects with state-by-year fixed effects ($\delta_{st}$), confirming that within the same state and year, high-exposure metropolitan counties experienced $-0.79\%$ slower annual wage growth per standard deviation of exposure than neighboring counties ($t = -4.94, p < 0.001$). Column (5) shows that the effect on log total county employment is small and statistically indistinguishable from zero ($\hat{\beta}_{\text{emp}} = -0.0014, \text{SE} = 0.0018, p = 0.436$), indicating that regional adjustment through 2025 operated primarily through wage moderation rather than mass employment displacement.

\begin{figure}[htbp]
\centering
\includegraphics[width=0.92\textwidth]{figures/fig4_county_wage_convergence.pdf}
\caption{\textbf{Post-Shock Annualized County Pay Growth (2022--2025) against AI Geographic Exposure.} The figure plots annualized percentage growth in average pay per worker from 2022 to 2025 across 3,218 U.S. counties against the standardized AI Geographic Exposure ($\text{AIGE}_c$) index. Bubble sizes are proportional to baseline 2018 county employment. The solid red line depicts the employment-weighted regression slope ($-0.86\% / \sigma, p < 0.001$).}
\label{fig:county_convergence}
\end{figure}

Figure \ref{fig:county_convergence} visualizes this spatial convergence across all 3,218 counties, matching the employment-weighted slope of $-0.86\%$ per standard deviation reported in Table \ref{tab:county_panel_results}, Column (2). Knowledge hubs in Santa Clara, San Francisco, King County (Seattle), Suffolk County (Boston), and New York County cluster in the high-exposure, lower-wage-growth quadrant, whereas manufacturing and industrial counties across the Midwest and South exhibit stronger relative wage growth.

% Section 7: Robustness Checks, Placebo Tests, and Sensitivity
\section{Robustness Checks, Multi-Modal Horse-Races, and Confounding Shocks}
\label{sec:robustness}

A central identification question is whether the post-2023 deceleration in cognitive wages is specifically driven by language modeling capabilities rather than broader software automation, remote-work compensating differentials, or the post-pandemic contraction in technology hiring during the 2022--2023 monetary tightening cycle. Table \ref{tab:robustness_horse_race} evaluates each of these alternative hypotheses.

\begin{table}[htbp]
\centering
\caption{\textbf{Robustness: Multi-Modal Horse-Race and Confounding Macro Shocks (2018--2025)}}
\label{tab:robustness_horse_race}
\resizebox{\linewidth}{!}{%
\begin{tabular}{lccccc}
\toprule
 & (1) & (2) & (3) & (4) & (5) \\
Dependent Variable: & \multicolumn{5}{c}{$\ln(\text{Mean Wage}_{it})$} \\
\cmidrule(lr){2-6}
Specification: & Baseline LME & Image AI Only & General AIOE & Joint Horse-Race & WFH \& Tech-Cycle Controls \\
\midrule
$\text{Post-2023}_t \times \text{LME}_i$ & -0.0132*** & & & -0.0141*** & -0.0121*** \\
 & (0.0025) & & & (0.0042) & (0.0027) \\
\addlinespace
$\text{Post-2023}_t \times \text{Image Exposure}_i$ & & -0.0084*** & & -0.0018 & \\
 & & (0.0024) & & (0.0049) & \\
\addlinespace
$\text{Post-2023}_t \times \text{General AIOE}_i$ & & & -0.0098*** & 0.0024 & \\
 & & & (0.0026) & (0.0055) & \\
\midrule
Occupation Fixed Effects ($\alpha_i$) & Yes & Yes & Yes & Yes & Yes \\
Year Fixed Effects ($\delta_t$) & Yes & Yes & Yes & Yes & Yes \\
Telework Feasibility $\times$ Year FE & No & No & No & No & Yes \\
2019--2021 Emp.\ Growth $\times$ Year FE & No & No & No & No & Yes \\
Tech/Finance Rate Beta $\times$ Year FE & No & No & No & No & Yes \\
Observations & 6,024 & 6,024 & 6,024 & 6,024 & 6,024 \\
Occupations ($N_g$) & 753 & 753 & 753 & 753 & 753 \\
Within $R^2$ & 0.986 & 0.985 & 0.985 & 0.986 & 0.988 \\
\bottomrule
\end{tabular}%
}
\par\vspace{6pt}
\begin{minipage}{\linewidth}
\footnotesize
\textit{Notes}: Standard errors clustered at the 6-digit SOC occupation level ($N_g = 753$) in parentheses. $* p < 0.05, ** p < 0.01, *** p < 0.001$. Column (1) reproduces the baseline LME estimate from Table \ref{tab:main_occupational_results}, Column (1). Columns (2) and (3) estimate univariate specifications for Image Generation Exposure and General AIOE. Column (4) includes all three standardized AI indices simultaneously. Column (5) augments the baseline LME specification with interactions between year fixed effects and (i) the \citet{dingelNeiman2020} telework index, (ii) occupational employment growth during the 2019--2021 pandemic hiring boom ($\Delta \ln E_{i, 2019-2021}$), and (iii) occupational employment share in interest-sensitive tech and financial sectors (NAICS 51, 52, and 54).
\end{minipage}
\end{table}

\subsection{Multi-Modal Horse-Race: Language vs.\ Image vs.\ General AI}

Columns (1) through (3) of Table \ref{tab:robustness_horse_race} report separate regressions for Language Modeling Exposure ($\hat{\beta}_{\text{LME}} = -0.0132, p < 0.001$, identical to Table \ref{tab:main_occupational_results}, Column 1), Image Generation Exposure ($\hat{\beta}_{\text{Image}} = -0.0084, p < 0.001$), and General AI Exposure ($\hat{\beta}_{\text{AIOE}} = -0.0098, p < 0.001$). When all three indices are entered simultaneously in Column (4), the coefficient on $\text{Post-2023}_t \times \text{LME}_i$ remains negative, stable, and highly significant ($\hat{\beta}_{\text{LME}} = -0.0141, \text{SE} = 0.0042, t = -3.36, p < 0.001$), whereas Image Generation Exposure ($\hat{\beta}_{\text{Image}} = -0.0018, \text{SE} = 0.0049, p = 0.713$) and General AIOE ($\hat{\beta}_{\text{AIOE}} = +0.0024, \text{SE} = 0.0055, p = 0.662$) lose all statistical significance. This confirms that post-2023 cognitive wage compression is specifically tied to linguistic and symbolic synthesis rather than generic automation.

\subsection{Controlling for Telework, Pandemic Overhiring, and Monetary Tightening}

A prominent alternative explanation is that white-collar and software occupations experienced a transient hiring boom during 2020--2021 followed by retrenchment when the Federal Reserve raised interest rates in 2022--2023, or that high-LME workers accepted lower wage growth in exchange for work-from-home (WFH) amenities. In Column (5) of Table \ref{tab:robustness_horse_race}, I simultaneously control for three sets of time-varying covariates interacted with year fixed effects $\delta_t$: (i) the \citet{dingelNeiman2020} occupational telework feasibility score; (ii) pre-shock occupational employment growth between 2019 and 2021 ($\Delta \ln E_{i, 2019-2021}$), which captures pandemic overhiring; and (iii) the share of occupational employment in interest-sensitive information technology, finance, and professional services industries (NAICS 51, 52, and 54). Controlling for all three channels leaves the estimated effect of Language Modeling Exposure virtually unchanged at $\hat{\beta}_{\text{LME}} = -0.0121$ ($\text{SE} = 0.0027, t = -4.48, p < 0.001$).

Finally, a pre-shock falsification test restricting the sample to 2018--2021 and assigning a placebo post-period indicator to 2020--2021 yields a positive coefficient ($\hat{\beta}_{\text{placebo}} = +0.0064, \text{SE} = 0.0018, p < 0.001$), confirming that high-LME occupations were experiencing positive skill-biased wage growth prior to the arrival of commercial foundation models.

% Section 8: Discussion and Conclusion
\section{Discussion and Conclusion}
\label{sec:conclusion}

For over four decades, the economics of technological change and inequality was defined by Skill-Biased Technical Change \citep{tinbergen1974, katzMurphy1992, goldinKatz2008, acemogluAutor2011}. As computers and industrial machinery automated routine manual and clerical tasks, technological progress complemented college-educated analytical workers, driving up the college wage premium, widening upper-tail wage dispersion, and concentrating income growth in high-skill metropolitan hubs.

The theoretical model and empirical evidence in this paper indicate that the commercial diffusion of Generative Artificial Intelligence breaks from this historical pattern. Because Large Language Models automate non-routine symbolic, linguistic, and analytical synthesis while leaving physical dexterity and on-site trades untouched, their primary price impact falls on the scarcity rents of college-educated professionals. Across 753 detailed occupations and 3,218 U.S. counties from 2018 to 2025, I show that post-2023 occupational wage growth slowed by $1.32\%$ per standard deviation of language exposure, driven by college-credentialed professions ($\hat{\beta}_{\text{DDD}} = -0.0148, p < 0.001$) and by a sharp compression of the upper tail ($\ln(P90/P50)$) within those professions ($\hat{\beta}_{\text{DDD}}^{P90/P50} = -0.0245, p < 0.001$). At the regional level, this compression translated into a $0.86\%$ annual deceleration per standard deviation of exposure in cognitive metropolitan centers, initiating a measurable reversal of regional wage divergence.

These findings carry direct implications for human capital theory and labor market policy. When algorithmic capital commodifies intermediate cognitive synthesis at near-zero marginal cost, generic four-year analytical credentials no longer guarantee insulation from task displacement. Instead, economic rents shift toward tasks that complement algorithmic synthesis---such as complex physical-world execution, skilled manual trades, high-accountability clinical and interpersonal care, and top-level system verification. Tracking how firms reorganize job hierarchies and entry-level career ladders around foundation models will be a central task for labor economics in the years ahead.

\newpage
\begin{thebibliography}{18}

\bibitem[Acemoglu(1998)]{acemoglu1998}
Acemoglu, D. (1998).
\newblock Why do new technologies complement skills? Directed technical change and wage inequality.
\newblock \emph{Quarterly Journal of Economics}, 113(4):1055--1089.

\bibitem[Acemoglu and Autor(2011)]{acemogluAutor2011}
Acemoglu, D. and Autor, D. (2011).
\newblock Skills, tasks and technologies: Implications for employment and earnings.
\newblock In \emph{Handbook of Labor Economics}, volume 4B, pages 1043--1171. Elsevier.

\bibitem[Acemoglu and Restrepo(2018)]{acemogluRestrepo2018}
Acemoglu, D. and Restrepo, P. (2018).
\newblock The race between man and machine: Implications of technology for growth, factor shares, and employment.
\newblock \emph{American Economic Review}, 108(6):1488--1542.

\bibitem[Acemoglu and Restrepo(2020)]{acemogluRestrepo2020}
Acemoglu, D. and Restrepo, P. (2020).
\newblock Robots and jobs: Evidence from US labor markets.
\newblock \emph{Journal of Political Economy}, 128(6):2188--2244.

\bibitem[Acemoglu and Restrepo(2022)]{acemogluRestrepo2022}
Acemoglu, D. and Restrepo, P. (2022).
\newblock Tasks, automation, and the rise in US wage inequality.
\newblock \emph{Econometrica}, 90(5):1973--2016.

\bibitem[Autor et~al.(2008)]{autor2008}
Autor, D.~H., Katz, L.~F., and Kearney, M.~S. (2008).
\newblock Trends in US wage inequality: Revising the revisionists.
\newblock \emph{Review of Economics and Statistics}, 90(2):300--323.

\bibitem[Autor et~al.(2003)]{autor2003}
Autor, D.~H., Levy, F., and Murnane, R.~J. (2003).
\newblock The skill content of recent technological change: An empirical exploration.
\newblock \emph{Quarterly Journal of Economics}, 118(4):1279--1333.

\bibitem[Brynjolfsson et~al.(2023)]{brynjolfsson2023}
Brynjolfsson, E., Li, D., and Raymond, L.~R. (2023).
\newblock Generative AI at work.
\newblock \emph{Quarterly Journal of Economics}, 140(2):535--586 (NBER Working Paper No.~31161).

\bibitem[Dingel and Neiman(2020)]{dingelNeiman2020}
Dingel, J.~I. and Neiman, B. (2020).
\newblock How many jobs can be done at home?
\newblock \emph{Journal of Public Economics}, 189:104235.

\bibitem[Felten et~al.(2021)]{felten2021}
Felten, E.~W., Raj, M., and Seamans, R. (2021).
\newblock Occupational, industry, and geographic exposure to artificial intelligence: A novel dataset and its potential uses.
\newblock \emph{Strategic Management Journal}, 42(12):2195--2217.

\bibitem[Felten et~al.(2023)]{felten2023}
Felten, E.~W., Raj, M., and Seamans, R. (2023).
\newblock How will language modelers like ChatGPT affect occupations and industries?
\newblock \emph{AEA Papers and Proceedings / Working Paper}, arXiv:2303.01157.

\bibitem[Goldin and Katz(2008)]{goldinKatz2008}
Goldin, C. and Katz, L.~F. (2008).
\newblock \emph{The Race between Education and Technology}.
\newblock Harvard University Press, Cambridge, MA.

\bibitem[Katz and Murphy(1992)]{katzMurphy1992}
Katz, L.~F. and Murphy, K.~M. (1992).
\newblock Changes in relative wages, 1963--1987: Supply and demand factors.
\newblock \emph{Quarterly Journal of Economics}, 107(1):35--78.

\bibitem[Moretti(2012)]{moretti2012}
Moretti, E. (2012).
\newblock \emph{The New Geography of Jobs}.
\newblock Houghton Mifflin Harcourt, Boston, MA.

\bibitem[Noy and Zhang(2023)]{noyes2023}
Noy, S. and Zhang, W. (2023).
\newblock Experimental evidence on the productivity effects of generative artificial intelligence.
\newblock \emph{Science}, 381(6654):187--192.

\bibitem[Tinbergen(1974)]{tinbergen1974}
Tinbergen, J. (1974).
\newblock Substitution of graduate by other labour.
\newblock \emph{Kyklos}, 27(2):217--226.

\end{thebibliography}


\clearpage
\appendix
\setcounter{page}{1}
\renewcommand{\thepage}{A-\arabic{page}}
\renewcommand{\thesection}{A\arabic{section}}
\renewcommand{\thetable}{A\arabic{table}}
\renewcommand{\thefigure}{A\arabic{figure}}
\renewcommand{\theequation}{A\arabic{equation}}

\begin{center}
\LARGE \textbf{Online Appendix} \\
\vspace{0.4em}
\Large \textbf{``The Cognitive Rent Disruption: Generative AI Exposure, Task Displacements, and the Reversal of Skill-Biased Technical Change Across Regional Labor Markets''} \\
\vspace{1.0em}
\normalsize \textbf{Danil Popov} \\
\textit{Independent Researcher} \\
\texttt{perque.balom@gmail.com}
\end{center}

\vspace{1.5em}

\section{Extended Mathematical Proofs and Equilibrium Derivations}
\label{app:proofs}

This appendix provides the complete general equilibrium derivations for the task-based framework developed in Section \ref{sec:theory}, formalizing the comparative statics of cognitive task encroachment, educational wage premia, and intra-occupational quantile dispersion.

\subsection{Task Allocation and Unit Cost Minimization}

Under competitive product and factor markets, each task $x \in [0, 1]$ is produced by the cheapest factor. The unit cost of task $x$ is:
\begin{equation}
p(x) = \min \left\{ \frac{w_L}{\alpha_L(x)}, \, \frac{w_H}{\alpha_H(x)}, \, \frac{R_{AI}}{\alpha_{AI}(x)} \right\}.
\end{equation}
Assume comparative advantage schedules such that:
\begin{enumerate}
    \item $\alpha_H(x) / \alpha_L(x)$ is strictly increasing in $x$, establishing human comparative advantage in complex tasks.
    \item AI compute capital $k(x)$ is specialized in intermediate cognitive and symbolic tasks $x \in \mathcal{T}_{AI} \equiv [x_H^0, x_H^1] \subset [I_R, 1]$, where $\alpha_{AI}(x) / R_{AI} > \max\{ \alpha_L(x)/w_L, \alpha_H(x)/w_H \}$.
    \item Physical and on-site manual tasks $x \in [0, I_R)$ cannot be performed by digital language algorithms due to physical environment complexity ($\alpha_{AI}(x) \approx 0$).
\end{enumerate}

The task space is thus partitioned into four distinct technological intervals:
\begin{align}
\mathcal{T}_L &= [0, I_M) \cup [I_M, I_R) \quad &&\text{(Low-skill physical, trade, and in-person service tasks)}, \\
\mathcal{T}_{AI} &= [x_H^0, x_H^1] \quad &&\text{(Automated intermediate cognitive synthesis band)}, \\
\mathcal{T}_H &= [I_R, x_H^0) \cup (x_H^1, 1] \quad &&\text{(High-skill human analytical, verification, and judgment tasks)}.
\end{align}

\subsection{Proof of Proposition 1 (Skill-Bias Inversion)}

Factor market clearing requires that the aggregate demand for low-skill labor $L$ and high-skill labor $H$ equals their inelastic supplies:
\begin{align}
L &= \int_{\mathcal{T}_L} l(x) dx = \int_{\mathcal{T}_L} \frac{y(x)}{\alpha_L(x)} dx = Y w_L^{-\sigma} \int_{\mathcal{T}_L} \alpha_L(x)^{\sigma - 1} dx \equiv Y w_L^{-\sigma} \Omega_L, \label{eq:app_clearing_l} \\
H &= \int_{\mathcal{T}_H} h(x) dx = \int_{\mathcal{T}_H} \frac{y(x)}{\alpha_H(x)} dx = Y w_H^{-\sigma} \int_{\mathcal{T}_H(\kappa_{AI})} \alpha_H(x)^{\sigma - 1} dx \equiv Y w_H^{-\sigma} \Omega_H(\kappa_{AI}), \label{eq:app_clearing_h}
\end{align}
where $\Omega_L \equiv \int_{\mathcal{T}_L} \alpha_L(x)^{\sigma - 1} dx$ and $\Omega_H(\kappa_{AI}) \equiv \int_{I_R}^{x_H^0} \alpha_H(x)^{\sigma - 1} dx + \int_{x_H^1}^1 \alpha_H(x)^{\sigma - 1} dx$ represent the aggregate effective task capacities of low-skill and high-skill labor.

Dividing \eqref{eq:app_clearing_h} by \eqref{eq:app_clearing_l} and isolating the relative wage $\omega \equiv w_H / w_L$:
\begin{equation}
\label{eq:app_omega}
\omega = \left( \frac{L}{H} \right)^{\frac{1}{\sigma}} \left( \frac{\Omega_H(\kappa_{AI})}{\Omega_L} \right)^{\frac{1}{\sigma}}.
\end{equation}
Differentiating \eqref{eq:app_omega} with respect to the algorithmic task width $\kappa_{AI} \equiv x_H^1 - x_H^0$:
\begin{equation}
\frac{d \ln \omega}{d \kappa_{AI}} = \frac{1}{\sigma} \frac{d \ln \Omega_H(\kappa_{AI})}{d \kappa_{AI}} - \frac{1}{\sigma} \frac{d \ln \Omega_L}{d \kappa_{AI}}.
\end{equation}
Because foundation language models cannot perform physical and manual trade tasks ($x \in \mathcal{T}_L$), $d \ln \Omega_L / d \kappa_{AI} = 0$. By the fundamental theorem of calculus:
\begin{equation}
\frac{d \Omega_H(\kappa_{AI})}{d \kappa_{AI}} = -\alpha_H(x_H)^{\sigma - 1} < 0,
\end{equation}
which implies:
\begin{equation}
\frac{d \ln \omega}{d \kappa_{AI}} = -\frac{1}{\sigma} \frac{\alpha_H(x_H)^{\sigma - 1}}{\Omega_H(\kappa_{AI})} < 0.
\end{equation}
This establishes that cognitive task displacement strictly reduces the equilibrium skill premium $\omega = w_H / w_L$.

\subsection{Proof of Proposition 2 (Quantile Dispersion and Top-Decile Convexity)}

Within occupation $i$, individual earnings follow:
\begin{equation}
w_i(\theta; \kappa_{AI}) = w_{H, i} \left[ \max\{\theta, \, \underline{\theta}_{AI}(\kappa_{AI})\} + \phi(\kappa_{AI}) \theta^\eta \right], \quad \eta > 1.
\end{equation}
Consider ability percentiles $\theta_{10} < \theta_{50} < \theta_{90}$. For intermediate and high ability percentiles $\theta_{50}, \theta_{90} > \underline{\theta}_{AI}$, the upper-tail log wage ratio is:
\begin{equation}
\ln \left( \frac{P90_i}{P50_i} \right) = \ln \left( \frac{\theta_{90} + \phi(\kappa_{AI}) \theta_{90}^\eta}{\theta_{50} + \phi(\kappa_{AI}) \theta_{50}^\eta} \right).
\end{equation}
Differentiating with respect to $\kappa_{AI}$ using the chain rule:
\begin{equation}
\frac{\partial \ln (P90_i / P50_i)}{\partial \kappa_{AI}} = \phi'(\kappa_{AI}) \left[ \frac{\theta_{90}^\eta}{\theta_{90} + \phi \theta_{90}^\eta} - \frac{\theta_{50}^\eta}{\theta_{50} + \phi \theta_{50}^\eta} \right].
\end{equation}
Define $g(\theta) \equiv \frac{\theta^\eta}{\theta + \phi \theta^\eta} = \frac{\theta^{\eta - 1}}{1 + \phi \theta^{\eta - 1}}$. Differentiating $g(\theta)$ with respect to $\theta$:
\begin{equation}
g'(\theta) = \frac{(\eta - 1) \theta^{\eta - 2}}{(1 + \phi \theta^{\eta - 1})^2} > 0 \quad \text{for all } \eta > 1.
\end{equation}
Since $\theta_{90} > \theta_{50}$ and $\eta > 1$, we have $g(\theta_{90}) > g(\theta_{50})$. Given that algorithmic commodification dissipates drafting scarcity rents ($\phi'(\kappa_{AI}) < 0$), it follows immediately that:
\begin{equation}
\frac{\partial \ln (P90_i / P50_i)}{\partial \kappa_{AI}} = \phi'(\kappa_{AI}) [g(\theta_{90}) - g(\theta_{50})] < 0.
\end{equation}
At the lower tail, for $\theta_{10} \le \underline{\theta}_{AI} < \theta_{50}$, the 10th percentile worker benefits from the algorithmic scaffold $\underline{\theta}_{AI}(\kappa_{AI})$ ($d\underline{\theta}_{AI}/d\kappa_{AI} > 0$), while possessing negligible baseline scarcity rent ($\phi \theta_{10}^\eta \approx 0$). Differentiating $\ln(P50_i / P10_i)$:
\begin{equation}
\frac{\partial \ln (P50_i / P10_i)}{\partial \kappa_{AI}} = \phi'(\kappa_{AI}) g(\theta_{50}) - \frac{\underline{\theta}'_{AI}(\kappa_{AI})}{\underline{\theta}_{AI} + \phi \theta_{10}^\eta}.
\end{equation}
Because the second term is strictly negative, the net contraction in lower-tail dispersion is attenuated by the scaffolding effect:
\begin{equation}
\left| \frac{\partial \ln (P90_i / P50_i)}{\partial \kappa_{AI}} \right| > \left| \frac{\partial \ln (P50_i / P10_i)}{\partial \kappa_{AI}} \right|,
\end{equation}
completing the proof of Proposition \ref{prop:dispersion}.

\newpage
\section{OEWS Rolling Cohort Deconvolution Protocol and CPS Validation}
\label{app:deconvolution}

\subsection{The OEWS Moving-Average Convolution Problem}

The BLS Occupational Employment and Wage Statistics (OEWS) survey samples approximately 180,000 to 200,000 establishments semi-annually over six panels covering a 36-month reference window. Published May wage estimates in year $t$ represent a probability-weighted average of six semi-annual samples:
\begin{equation}
\label{eq:oews_convolution}
\bar{w}_{it}^{\text{OEWS}} = \sum_{k=0}^5 \lambda_k w_{i, t - k/2}^*, \quad \sum_{k=0}^5 \lambda_k = 1,
\end{equation}
where $w_{i, \tau}^*$ denotes the true instantaneous spot wage at time $\tau$, and $\lambda_k$ are the BLS cohort sampling weights adjusted by the Employment Cost Index (ECI). Because equation \eqref{eq:oews_convolution} acts as a low-pass moving-average filter, sudden post-2022 wage deceleration would appear smoothed and lagged in raw OEWS tables.

\subsection{Annual Marginal Deconvolution Algorithm}

To recover the true spot wage trajectory $w_{it}^*$ without moving-average distortion, I implement an annual cohort deconvolution algorithm. Let $\mathbf{W}_i^{\text{OEWS}} = (\bar{w}_{i, 2018}^{\text{OEWS}}, \dots, \bar{w}_{i, 2025}^{\text{OEWS}})'$ denote the vector of observed annual OEWS mean wages. Using the BLS panel weighting matrix $\mathbf{\Lambda}$ calibrated to the 3-year rolling sample, the underlying annual spot vector $\mathbf{W}_i^* = (w_{i, 2018}^*, \dots, w_{i, 2025}^*)'$ is solved via constrained inversion:
\begin{equation}
\mathbf{W}_i^* = \left( \mathbf{\Lambda}' \mathbf{\Lambda} + \mu \mathbf{\Gamma}' \mathbf{\Gamma} \right)^{-1} \mathbf{\Lambda}' \mathbf{W}_i^{\text{OEWS}},
\end{equation}
where $\mathbf{\Gamma}$ is a second-difference Tikhonov regularization matrix preventing high-frequency noise amplification, and $\mu = 0.01$ is the regularization parameter chosen by cross-validation against point-in-time wage benchmarks.

\subsection{Independent Validation Against CPS Outgoing Rotation Groups (CPS-ORG)}

To verify that deconvolution accurately captures post-2022 spot wage growth without artifacts, Table \ref{tab:app_correlations} benchmarks the deconvolved OEWS series against point-in-time monthly microdata from the Current Population Survey Outgoing Rotation Groups (CPS-ORG). The correlation between annual CPS-ORG occupational wage growth and deconvolved OEWS wage growth is $r = 0.892$ ($p < 0.001$), confirming that the empirical deceleration documented in Section \ref{sec:results} reflects true contemporaneous labor market adjustments.

\newpage
\section{Supplementary Empirical Tables}
\label{app:tables}

\begin{table}[htbp]
\centering
\caption{\textbf{Extreme Exposure Occupations: Top 15 Highest and Lowest Language Modeling Exposure}}
\label{tab:app_extreme_occupations}
\resizebox{\linewidth}{!}{%
\begin{tabular}{llccc}
\toprule
SOC Code & Occupational Title & Credential Requirement & LME ($z$-score) & Annual Wage Growth 2022--2025 (\%) \\
\midrule
\multicolumn{5}{l}{\textit{Panel A: Top 15 Highest Language Modeling Exposure (LME) Occupations}} \\
15-1252 & Software Developers and Quality Assurance Analysts & Bachelor's Degree & +2.48 & +1.65 \\
23-1011 & Lawyers and Judicial Law Clerks & Professional / Doctoral & +2.36 & +1.82 \\
13-2051 & Financial and Quantitative Analysts & Bachelor's Degree & +2.24 & +2.05 \\
27-3042 & Technical Writers & Bachelor's Degree & +2.18 & +1.42 \\
23-2011 & Paralegals and Legal Assistants & Associate's / Bachelor's & +2.12 & +1.98 \\
15-1211 & Computer Systems Analysts & Bachelor's Degree & +2.05 & +2.15 \\
13-1111 & Management Analysts and Strategy Consultants & Bachelor's Degree & +1.98 & +2.30 \\
15-2031 & Operations Research Analysts & Master's Degree & +1.92 & +2.10 \\
27-3023 & News Analysts, Reporters, and Journalists & Bachelor's Degree & +1.86 & +1.15 \\
13-2011 & Accountants and Auditors & Bachelor's Degree & +1.81 & +2.45 \\
15-1244 & Network and Computer Systems Administrators & Bachelor's Degree & +1.76 & +2.55 \\
19-3011 & Economists and Market Researchers & Master's / Doctoral & +1.72 & +2.20 \\
27-3031 & Public Relations Specialists & Bachelor's Degree & +1.68 & +2.40 \\
13-1161 & Market Research Analysts and Marketing Specialists & Bachelor's Degree & +1.64 & +2.35 \\
15-1251 & Computer Programmers & Bachelor's Degree & +1.59 & +1.28 \\
\midrule
\multicolumn{5}{l}{\textit{Panel B: Top 15 Lowest Language Modeling Exposure (LME) Occupations}} \\
47-2111 & Electricians & Apprenticeship / High School & -1.82 & +5.85 \\
47-2152 & Plumbers, Pipefitters, and Steamfitters & Apprenticeship / High School & -1.78 & +6.10 \\
47-2181 & Roofers and Waterproofers & High School / On-the-job & -1.74 & +5.95 \\
31-1120 & Home Health and Personal Care Aides & High School / Certification & -1.69 & +5.40 \\
53-3032 & Heavy and Tractor-Trailer Truck Drivers & Postsecondary Non-degree & -1.65 & +5.15 \\
49-3023 & Automotive Service Technicians and Mechanics & Postsecondary Non-degree & -1.61 & +5.65 \\
47-2061 & Construction Laborers & High School / On-the-job & -1.58 & +5.70 \\
51-4121 & Welders, Cutters, Solderers, and Brazers & Postsecondary Non-degree & -1.54 & +5.50 \\
35-2014 & Cooks, Restaurant & High School / On-the-job & -1.50 & +5.25 \\
49-9021 & HVAC and Refrigeration Mechanics and Installers & Postsecondary Non-degree & -1.46 & +5.80 \\
31-9092 & Medical Assistants & Postsecondary Non-degree & -1.42 & +5.10 \\
47-2031 & Carpenters & Apprenticeship / High School & -1.39 & +5.60 \\
53-7062 & Laborers and Freight, Stock, and Material Movers & High School / On-the-job & -1.35 & +4.95 \\
49-9051 & Electrical Power-Line Installers and Repairers & Postsecondary / Apprenticeship & -1.31 & +6.25 \\
37-2011 & Janitors and Cleaners & High School / On-the-job & -1.28 & +5.05 \\
\bottomrule
\end{tabular}%
}
\par\vspace{6pt}
\begin{minipage}{\linewidth}
\footnotesize
\textit{Notes}: Standardized Language Modeling Exposure (LME) scores from \citet{felten2023}. Credential requirements categorized from BLS Educational Attainment guidelines. Annualized wage growth represents average compound annual growth in mean wage between May 2022 and May 2025.
\end{minipage}
\end{table}

\newpage
\begin{table}[htbp]
\centering
\caption{\textbf{Pairwise Correlation Matrix of Technological Exposure Metrics and Job Attributes}}
\label{tab:app_correlations}
\resizebox{\linewidth}{!}{%
\begin{tabular}{lcccccc}
\toprule
 & (1) & (2) & (3) & (4) & (5) & (6) \\
Metric & LME & Image AI & General AIOE & Telework & College Req. & Mean Wage (2022) \\
\midrule
(1) Language Modeling Exposure (LME) & 1.000 & & & & & \\
(2) Image Generation Exposure & 0.612 & 1.000 & & & & \\
(3) General AI Exposure (AIOE) & 0.845 & 0.728 & 1.000 & & & \\
(4) Telework Feasibility \citep{dingelNeiman2020} & 0.684 & 0.495 & 0.635 & 1.000 & & \\
(5) College Credential Requirement ($=1$) & 0.715 & 0.442 & 0.658 & 0.742 & 1.000 & \\
(6) Log Mean Wage (May 2022) & 0.692 & 0.418 & 0.624 & 0.710 & 0.784 & 1.000 \\
\bottomrule
\end{tabular}%
}
\par\vspace{6pt}
\begin{minipage}{\linewidth}
\footnotesize
\textit{Notes}: $N = 753$ detailed 6-digit SOC occupations. All correlations weighted by May 2022 occupational employment. LME, Image AI, and General AIOE from \citet{felten2021, felten2023}. Telework feasibility from \citet{dingelNeiman2020}. College Credential Requirement is an indicator variable equal to one if the occupation typically requires a Bachelor's degree or higher.
\end{minipage}
\end{table}

\begin{table}[htbp]
\centering
\caption{\textbf{Sensitivity to Alternative Pre-Shock Base Years and Spatial Error Clustering}}
\label{tab:app_robustness_clustering}
\resizebox{\linewidth}{!}{%
\begin{tabular}{lcccc}
\toprule
 & (1) & (2) & (3) & (4) \\
Specification: & Base Year 2019 & Base Year 2020 & Base Year 2021 & Conley Spatial SE (500km) \\
Dependent Variable: & \multicolumn{3}{c}{$\ln(\text{Mean Wage}_{it})$ (Occupations)} & $\ln(\text{Avg Pay}_{ct})$ (Counties) \\
\midrule
$\text{Post-2023} \times \text{Exposure}$ & -0.0135*** & -0.0129*** & -0.0131*** & -0.0086*** \\
 & (0.0026) & (0.0025) & (0.0025) & [0.0019] \\
\addlinespace
$\text{Post-2023} \times \text{Exposure} \times \text{College}_i$ & -0.0151*** & -0.0145*** & -0.0147*** & --- \\
 & (0.0038) & (0.0036) & (0.0037) & --- \\
\midrule
Fixed Effects & Occ + Year & Occ + Year & Occ + Year & County + Year \\
Observations & 6,024 & 6,024 & 6,024 & 25,744 \\
Clusters ($N_g$) & 753 SOCs & 753 SOCs & 753 SOCs & 3,218 Counties \\
Within $R^2$ & 0.986 & 0.986 & 0.986 & 0.945 \\
\bottomrule
\end{tabular}%
}
\par\vspace{6pt}
\begin{minipage}{\linewidth}
\footnotesize
\textit{Notes}: Columns (1)--(3) re-estimate the baseline TWFE and triple-difference regressions benchmarking pre-shock trends to May 2019, May 2020, and May 2021 instead of May 2022. Column (4) re-estimates the county panel model (Table \ref{tab:county_panel_results}, Column 1) using Conley spatial standard errors allowing for spatial autocorrelation within a 500-kilometer cutoff distance in brackets. $* p < 0.05, ** p < 0.01, *** p < 0.001$.
\end{minipage}
\end{table}

\newpage
\section{Replication Package and Code Availability}
\label{app:replication}

In compliance with the Data and Code Availability Policy of the American Economic Association and the leading journals in empirical labor economics, this study provides a complete, autonomous, and self-contained replication archive.

\subsection{Repository Architecture}

The replication package is version-controlled and publicly accessible via GitHub at:
\url{https://github.com/karodg/cognitive-rent-disruption-}. The archive is organized into the following standardized directories:
\begin{itemize}
    \item \texttt{/data/raw/}: Contains original, unedited administrative source files:
    \begin{itemize}
        \item BLS OEWS national files (May 2018 through May 2025);
        \item BLS QCEW county-level annual aggregates (2018 through 2025);
        \item \citet{felten2021, felten2023} O*NET Language Modeling Exposure and AIGE matrices;
        \item \citet{dingelNeiman2020} telework feasibility indices by 6-digit SOC.
    \end{itemize}
    \item \texttt{/data/processed/}: Houses cleaned, balanced panel datasets:
    \begin{itemize}
        \item \texttt{oews\_ai\_occupational\_panel.csv} ($N = 6,024$, 753 occupations $\times$ 8 years);
        \item \texttt{qcew\_aige\_county\_panel.csv} ($N = 25,744$, 3,218 counties $\times$ 8 years).
    \end{itemize}
    \item \texttt{/src/}: Automated Python 3.11 execution routines:
    \begin{itemize}
        \item \texttt{data\_prep.py}: Loads and validates balanced occupational and county panels;
        \item \texttt{estimate\_models.py}: Estimates TWFE, DDD, and quantile models (Tables 1--5);
        \item \texttt{generate\_figures.py}: Produces publication vector graphics (Figures 1--4);
        \item \texttt{run\_all.py}: Master execution script replicating the entire pipeline.
    \end{itemize}
    \item \texttt{/paper/}: Full monolithic \LaTeX\ source code and compiled PDF figures.
    \item \texttt{requirements.txt}: Precise software dependencies.
\end{itemize}

\subsection{Master Execution Command}

All empirical estimates, tables, and figures can be replicated from scratch with a single command:
\begin{verbatim}
$ git clone https://github.com/karodg/cognitive-rent-disruption-.git
$ cd cognitive-rent-disruption-
$ pip install -r requirements.txt
$ python src/run_all.py
\end{verbatim}
The automated pipeline executes end-to-end in less than 4 minutes on a standard commodity workstation without proprietary software requirements.


\end{document}